\scanpage{4}
$C_{10}(S)$ Catégorie des torseurs sous $(\mathbb Z/2)_S\mathbin{\cdot}_{1/2}\mathrm G_{m\,S}$.

\deleted{\uline{\illegible{1}}} Toutes ces catégories sont canoniquemt équivalentes...

\deleted{Condition équi} Pour $S'$ sur $S$ fixé (NB $S'=\uline{\operatorname{Spec}}(A)\ldots$) les structures en question forment une \inserted{\unsure{Gr-|cat}} catégorie (clair sur $C_2$ et sur $\unsure{C_3|C_0}$). Si $T$ est un objet de celle-ci ($\unsure{\mathcal U|U}$-torseur), on veut que $T^{\wedge 2}$ est \uline{loc. trivial au sens de Zariski}.
% slanted
\margin{faux, m\textsuperscript{m} sur un corps p.ex. sur $\mathbb R$ la forme quadratique \unsure{$a$|$\alpha$}$(x^2+y^2)$ n'est pas isom. à $xy\ldots$}
 On veut expliciter ce qu'est une trivialisation de $T^{\wedge 2}$. Par la suite exacte de Kummer $0-\mu_2-\mathcal U\xrightarrow{2}\mathcal U\longrightarrow 0$, on voit qu'il revient au $\unsure{=}$ de se donner une restriction du gpe structural à $\mu_2$, ou une orbite \deleted{sous} de $T$ sous $\mu_2$, ou un morphisme de degr $2$ $X\longrightarrow X'$, où $X'$ est un fibré en droites avec diviseur \inserted{relatif} étale $\Delta'$, et une section \deleted{\unsure{de}} de $X'\setminus\Delta'$, $X-X'$ induisant un morphisme \inserted{fini} étale
% slanted
\margin{\uline{NB} L'action de $\mu_2$ sur $(X,\Delta)$ est l'unique action fidèle qui induit l'\uline{identité} sur $\Delta\ldots$ Cf \uline{définition} antérieure d'une cat.\ldots)}
 $X\setminus\Delta\longrightarrow X'\setminus\Delta'$. \unsure{Celà|Cel|Ceci} signifie aussi qu'on a une restriction du Groupe structural au ss-gpe $(\mathbb Z/2)_S\mathbin{\cdot}_{1/2}\mu_{2\,S}\simeq(\mathbb Z/2)_S\times(\mu_2)_S$ agissant sur $\underline O_S\times\underline O_S$ par symétrie et homothéties, Donc la catégorie de ces objets
% slanted
\margin{\uline{NB} $X'=X/\mu_2$ est défini canoniquement, le difficile est de \unsure{m\textsuperscript{e}|mettre} la section de $X'\setminus\Delta'$. \unsure{Chois.|Chos.} $\mathbb Z/2\mathbb Z=\mathbb D_S=\underline O(2)_S$ ($1/2$ \unsure{divisé|disons})}
 s'identifie à la cat. des triples $(A,L,\varepsilon)$, $A$ Algèbre finie étale de rg $2$, sur $S$, $L$ mod. inv. sur $S$, $\varepsilon$ une trivialisation de $L^{\otimes2}$ -- on prend $V=A\otimes_{\underline O_S}L$, d'où $N_{A/\underline O_S}V\simeq L^{\otimes2}$, et on prend $\xi=\varepsilon\ldots$

% archive sheet number: 3
